% =====================================================================
%  FUNDAMENTOS DE MECÂNICA — Apostila de Revisão com Personagens Pokémon
%  Profª Camila Leite Coura Mariano de Oliveira
%  Conteúdo: revisao.tex + revisao2.tex + movbim.tex
% =====================================================================

\documentclass[10.5pt, a4paper]{extarticle}

% ── Pacotes Fundamentais ─────────────────────────────────────────────
\usepackage[utf8]{inputenc}
\usepackage[T1]{fontenc}
\usepackage[brazil]{babel}
\usepackage{amsmath, amssymb, amsthm}
\usepackage{geometry}
\usepackage[table]{xcolor}
\usepackage{graphicx}
\usepackage{titlesec}
\usepackage{tcolorbox}
\usepackage{enumitem}
\usepackage{fancyhdr}
\usepackage{tikz}
\usepackage{pgfplots}
\pgfplotsset{compat=1.18}
\usepackage{hyperref}
\usepackage{fontawesome5}
\usepackage{setspace}
\usepackage{scrextend}
\usepackage{booktabs}

\usetikzlibrary{arrows.meta, calc}

% ── Configurações de Layout ──────────────────────────────────────────
\onehalfspacing
\changefontsizes{10.5pt}
\geometry{margin=2.5cm}

% ── Cores temáticas Pokémon ──────────────────────────────────────────
\definecolor{pokered}{RGB}{206,30,30}
\definecolor{pokeblue}{RGB}{58,90,174}
\definecolor{pokeyellow}{RGB}{250,208,0}
\definecolor{pokegray}{RGB}{100,100,100}
\definecolor{pokegreen}{RGB}{56,142,60}
\definecolor{pokeorange}{RGB}{230,100,20}
\definecolor{lightred}{RGB}{255,235,235}
\definecolor{lightblue}{RGB}{235,245,255}
\definecolor{lightyellow}{RGB}{255,252,230}
\definecolor{lightgreen}{RGB}{235,250,235}

% ── Cabeçalho / Rodapé ───────────────────────────────────────────────
\pagestyle{fancy}
\fancyhf{}
\fancyhead[L]{\textcolor{pokered}{\textbf{Fundamentos de Mecânica}}}
\fancyhead[R]{\textcolor{pokeblue}{\textbf{Notas de Aula}}}
\fancyfoot[C]{\thepage}
\renewcommand{\headrulewidth}{1.5pt}
\renewcommand{\headrule}{\hbox to\headwidth{\color{pokered}\leaders\hrule height \headrulewidth\hfill}}

% ── Estilo de Títulos ────────────────────────────────────────────────
\titleformat{\section}{\Large\bfseries\color{pokered}}{}{0em}{\faFireAlt\quad}[]
\titleformat{\subsection}{\large\bfseries\color{pokeblue}}{}{0em}{\faChevronRight\quad}[]
\titleformat{\subsubsection}{\normalsize\bfseries\color{pokegreen}}{}{0em}{}[]

% ── Caixas Temáticas (TColorBox) ─────────────────────────────────────
\tcbuselibrary{skins, breakable}

\newtcolorbox{pikabox}[1][]{
  enhanced, breakable,
  title={\faLightbulb\quad Pikachu diz:},
  colback=lightyellow, colframe=pokeyellow,
  colbacktitle=pokeyellow, coltitle=black,
  fonttitle=\bfseries, boxrule=1.5pt,
  attach boxed title to top left={yshift=-2mm,xshift=4mm},
  #1
}

\newtcolorbox{charmanderbox}[1][]{
  enhanced, breakable,
  title={\faBurn\quad Charmander Alerta:},
  colback=lightred, colframe=pokered,
  colbacktitle=pokered, coltitle=white,
  fonttitle=\bfseries, boxrule=1.5pt,
  attach boxed title to top left={yshift=-2mm,xshift=4mm},
  #1
}

\newtcolorbox{squirtlebox}[1][]{
  enhanced, breakable,
  title={\faTint\quad Squirtle Explica:},
  colback=lightblue, colframe=pokeblue,
  colbacktitle=pokeblue, coltitle=white,
  fonttitle=\bfseries, boxrule=1.5pt,
  attach boxed title to top left={yshift=-2mm,xshift=4mm},
  #1
}

\newtcolorbox{bulbabox}[1][]{
  enhanced, breakable,
  title={\faLeaf\quad Bulbasaur Resolve:},
  colback=lightgreen, colframe=pokegreen,
  colbacktitle=pokegreen, coltitle=white,
  fonttitle=\bfseries, boxrule=1.5pt,
  attach boxed title to top left={yshift=-2mm,xshift=4mm},
  #1
}

\newtcolorbox{definicaobox}[1][]{
  enhanced, breakable,
  title={\faBookOpen\quad Definição},
  colback=gray!5, colframe=pokegray,
  colbacktitle=pokegray, coltitle=white,
  fonttitle=\bfseries, boxrule=1pt,
  attach boxed title to top left={yshift=-2mm,xshift=4mm},
  #1
}

% ── Hiperlinks ───────────────────────────────────────────────────────
\hypersetup{colorlinks=true, linkcolor=pokered, urlcolor=pokeblue}

% =====================================================================
\begin{document}
% =====================================================================

% ── CAPA ─────────────────────────────────────────────────────────────
\begin{titlepage}
  \centering
  \vspace*{2cm}
  {\Huge\bfseries\color{pokered} FUNDAMENTOS DE MECÂNICA\par}
  \vspace{0.5cm}
  {\LARGE\color{pokered} Apostila de Revisão com Personagens Pokémon\par}
  \vspace{0.3cm}
  {\large\color{pokered} Profª Camila Leite Coura Mariano de Oliveira\par}
  \vspace{1.5cm}
  \rule{0.8\textwidth}{0.4pt}
  \vspace{1cm}

  % ── POKÉBOLA (estrutura da professora) ──────────────────────────────
  \begin{tikzpicture}
    % 1. Círculo principal (fundo branco + contorno preto)
    \draw[fill=white, line width=1.5pt] (0,0) circle (2cm);
    % 2. Metade superior vermelha
    \begin{scope}
      \clip (0,0) circle (2cm);
      \fill[pokered] (-2,0) rectangle (2,2);
    \end{scope}
    % 3. Faixa central preta
    \fill[black] (-2,-0.1) rectangle (2,0.1);
    % 4. Botão externo preto
    \draw[fill=black] (0,0) circle (0.6cm);
    % 5. Botão médio branco
    \draw[fill=white] (0,0) circle (0.4cm);
    % 6. Botão interno (menor)
    \draw[fill=white, line width=0.5pt] (0,0) circle (0.25cm);
    % 7. Contorno externo final
    \draw[line width=1.5pt] (0,0) circle (2cm);
  \end{tikzpicture}

  \vspace{1cm}
  \rule{0.8\textwidth}{0.4pt}
  \vspace{1cm}

  {\Large Com os personagens:\par}
  \vspace{0.4cm}
  {\large
    \faLightbulb\ \textbf{Pikachu} --- Dicas e curiosidades\\[4pt]
    \faBurn\ \textbf{Charmander} --- Alertas e armadilhas\\[4pt]
    \faTint\ \textbf{Squirtle} --- Explicações conceituais\\[4pt]
    \faLeaf\ \textbf{Bulbasaur} --- Exemplos resolvidos
  }

  \vfill
  {\normalsize Aulas de Mecânica Clássica $\cdot$ Cinemática e Dinâmica\par}
\end{titlepage}

\pagecolor{white}\color{black}

% ── SUMÁRIO ──────────────────────────────────────────────────────────

\tableofcontents
\newpage

% =====================================================================
\section{Movimento Unidimensional: Conceitos Fundamentais}
% =====================================================================

\subsection{Posição, Deslocamento e Distância}

\begin{definicaobox}
A \textbf{posição} $x(t)$ de uma partícula é sua localização ao longo de um eixo de referência em função do tempo $t$.

O \textbf{deslocamento} é a variação de posição:
\[
  \Delta x = x_f - x_i
\]
A \textbf{distância percorrida} é o comprimento total do caminho (sempre $\geq 0$).
\end{definicaobox}

\begin{pikabox}
Deslocamento pode ser \textbf{negativo}! Se você sai de $x=5\,\text{m}$ e vai para $x=2\,\text{m}$, o deslocamento é $\Delta x = -3\,\text{m}$, mas a distância percorrida é $3\,\text{m}$.
\end{pikabox}

\subsection{Velocidade Média}

A \textbf{velocidade média} é definida como:
\[
  \bar{v} = \frac{\Delta x}{\Delta t} = \frac{x_f - x_i}{t_f - t_i}
\]

\begin{charmanderbox}
\textbf{Não confunda} velocidade média com rapidez média!
\begin{itemize}[nosep]
  \item \textbf{Velocidade média}: grandeza com sinal, depende apenas dos pontos inicial e final.
  \item \textbf{Rapidez média}: distância total dividida pelo tempo total (sempre positiva).
\end{itemize}
\end{charmanderbox}

\subsection{Velocidade Instantânea e Derivada}

\begin{definicaobox}
A velocidade instantânea $v(t)$ é a taxa de variação da posição no tempo — o limite da velocidade média quando $\Delta t \to 0$:
\[
  v(t) = \lim_{\Delta t \to 0} \frac{\Delta x}{\Delta t} = \frac{dx}{dt}
\]
\end{definicaobox}

Geometricamente, $v(t)$ é a \textit{inclinação da tangente} ao gráfico $x(t)$.

\begin{squirtlebox}
Se o gráfico $x \times t$ for uma parábola, sua derivada (velocidade) será uma reta. Geometricamente, a velocidade em um ponto é a inclinação da reta tangente à curva naquele instante.
\end{squirtlebox}

\begin{center}
\begin{tikzpicture}
  \begin{axis}[
    width=10cm, height=5cm,
    xlabel={$t$ (s)}, ylabel={$x$ (m)},
    title={Gráfico Posição $\times$ Tempo},
    grid=major, grid style={dotted, pokegray!40},
    axis line style={pokegray},
    title style={color=pokered, font=\bfseries},
    xlabel style={color=pokegray}, ylabel style={color=pokegray},
    xmin=0, xmax=5, ymin=0, ymax=12
  ]
    \addplot[thick, pokeblue, samples=100, domain=0:5] {0.5*x^2 + 0.5};
    \addplot[thick, pokered, dashed, domain=1.5:4.5] {2*x - 1.5};
    \node[pokered, font=\small] at (axis cs:4.2, 7) {tangente em $t_0$};
  \end{axis}
\end{tikzpicture}
\end{center}

% =====================================================================
\section{O Problema Inverso: da Velocidade ao Deslocamento}
% =====================================================================

\subsection{Interpretação Gráfica}

\begin{squirtlebox}
Se conhecemos $v(t)$, podemos recuperar $\Delta x$ calculando a \textbf{área sob a curva} $v(t)$ no intervalo $[t_i, t_f]$:
\[
  \Delta x = \int_{t_i}^{t_f} v(t)\,dt
\]
Esta é a ideia central do \textit{Problema Inverso da Velocidade}.
\end{squirtlebox}

\subsection{Aceleração Média e Instantânea}

\begin{definicaobox}
\[
  \bar{a} = \frac{\Delta v}{\Delta t}, \qquad
  a(t) = \frac{dv}{dt} = \frac{d^2x}{dt^2}
\]
\end{definicaobox}

\begin{pikabox}
\textbf{Sinal da aceleração vs.\ sinal da velocidade:}
\begin{itemize}[nosep]
  \item $v > 0$ e $a > 0$: movimento acelerado no sentido positivo.
  \item $v > 0$ e $a < 0$: movimento desacelerado no sentido positivo.
  \item $v < 0$ e $a < 0$: movimento acelerado no sentido negativo.
  \item $v < 0$ e $a > 0$: movimento desacelerado no sentido negativo.
\end{itemize}
\end{pikabox}

% =====================================================================
\section{Cinemática Escalar Unidimensional: MRUV e Queda Livre}
% =====================================================================

\subsection{Movimento Retilíneo Uniforme (MRU)}

\begin{definicaobox}
No MRU, a aceleração é nula: $a = 0$, portanto $v = $ constante.
\[
  \boxed{x(t) = x_0 + v \cdot t}
\]
\end{definicaobox}

\begin{center}
\begin{tikzpicture}
  \begin{axis}[
    width=5.5cm, height=4cm,
    xlabel={$t$}, ylabel={$x$},
    title={$x \times t$ (MRU)},
    title style={color=pokered, font=\bfseries\small},
    xmin=0, xmax=5, ymin=0, ymax=8,
    grid=major, grid style={dotted, gray!40},
    axis lines=left
  ]
    \addplot[thick, pokeblue] {0.5 + 1.4*x};
  \end{axis}
\end{tikzpicture}
\hspace{1cm}
\begin{tikzpicture}
  \begin{axis}[
    width=5.5cm, height=4cm,
    xlabel={$t$}, ylabel={$v$},
    title={$v \times t$ (MRU)},
    title style={color=pokered, font=\bfseries\small},
    xmin=0, xmax=5, ymin=0, ymax=4,
    grid=major, grid style={dotted, gray!40},
    axis lines=left
  ]
    \addplot[thick, pokered] {1.4};
  \end{axis}
\end{tikzpicture}
\end{center}

\subsection{Encontro de Dois Móveis no MRU}

Dois móveis $A$ e $B$ com posições $x_A(t) = x_{0A} + v_A t$ e $x_B(t) = x_{0B} + v_B t$ se encontram quando $x_A = x_B$.

\begin{bulbabox}
\textbf{Exemplo:} $x_{0A}=0\,\text{m}$, $v_A=20\,\text{m/s}$; $x_{0B}=200\,\text{m}$, $v_B=-10\,\text{m/s}$.
\[
  20t = 200 - 10t \;\Rightarrow\; 30t = 200 \;\Rightarrow\; t = \frac{20}{3}\,\text{s} \approx 6{,}7\,\text{s}
\]
\[
  x_{\text{encontro}} = 20 \cdot \frac{20}{3} = \frac{400}{3}\,\text{m} \approx 133{,}3\,\text{m}
\]
\end{bulbabox}

\subsection{MRUV e a Equação de Torricelli}

Quando a aceleração é constante, temos o \textbf{Movimento Retilíneo Uniformemente Variado (MRUV)}:
\[
  v(t) = v_0 + at, \qquad x(t) = x_0 + v_0 t + \tfrac{1}{2}at^2
\]

\begin{pikabox}
Se você não tiver o tempo, use \textbf{Torricelli}! É a salvação para muitos problemas:
\[
  v^2 = v_0^2 + 2a\,\Delta x
\]
\end{pikabox}

\subsection{Queda Livre e Lançamento Vertical}

\begin{charmanderbox}
Na queda livre, o objeto é solto ($v_0 = 0$). No lançamento vertical, a aceleração $g \approx 9{,}8\,\text{m/s}^2$ sempre aponta para \textbf{baixo}, mesmo quando o objeto está subindo!
\end{charmanderbox}

\begin{bulbabox}
\textbf{Exemplo do Pidgey:} Um Pidgey solta uma Pokébola de uma altura de $20\,\text{m}$. Qual a velocidade ao atingir o solo? ($g=10\,\text{m/s}^2$)
\[
  v^2 = 0^2 + 2(10)(20) \;\Rightarrow\; v^2 = 400 \;\Rightarrow\; v = 20\,\text{m/s}
\]
\end{bulbabox}

% =====================================================================
\section{Movimentos Bidimensionais e Vetores}
% =====================================================================

\subsection{Introdução ao Movimento Bidimensional}

O movimento bidimensional descreve a trajetória de um corpo em um plano, exigindo o uso de vetores para determinar sua localização e evolução temporal.

\begin{center}
\begin{tikzpicture}[>=stealth]
  % Grade
  \draw[very thin, gray!30] (-1,-1) grid (6,6);
  % Eixos
  \draw[->, thick] (-1,0) -- (6,0) node[right] {$x$};
  \draw[->, thick] (0,-1) -- (0,6) node[above] {$y$};
  % Marcadores
  \foreach \x in {1,2,3,4,5}
    \draw (\x, 1pt) -- (\x, -1pt) node[below] {\footnotesize \x};
  \foreach \y in {1,2,3,4,5}
    \draw (1pt, \y) -- (-1pt, \y) node[left] {\footnotesize \y};
  % Vetor posição
  \draw[->, blue, ultra thick] (0,0) -- (4,4)
    node[midway, above left, black] {$\vec{r} = 4\hat{i} + 4\hat{j}$};
  % Projeções
  \draw[dashed, gray] (4,0) -- (4,4);
  \draw[dashed, gray] (0,4) -- (4,4);
  \filldraw[black] (4,4) circle (1.5pt) node[anchor=south west] {$(4,4)$};
  % Ângulo
  \draw[->, red] (1,0) arc (0:45:1);
  \node[red] at (1.3, 0.5) {$45^\circ$};
\end{tikzpicture}
\end{center}

\subsection{Grandezas Vetoriais e Escalares}

\begin{definicaobox}
\begin{tabular}{ll}
  \textbf{Escalares}: & Massa, temperatura, distância, rapidez. \\
  \textbf{Vetoriais}: & Posição, deslocamento, velocidade, aceleração, força.
\end{tabular}
\end{definicaobox}

\subsection{Vetor Posição}

O vetor posição $\vec{r}(t)$ localiza a partícula em relação à origem de um sistema de coordenadas:
\[
  \vec{r}(t) = x(t)\,\hat{i} + y(t)\,\hat{j}
\]

\subsection{Operações com Vetores}

Dado $\vec{A} = A_x\,\hat{i} + A_y\,\hat{j}$ e $\vec{B} = B_x\,\hat{i} + B_y\,\hat{j}$:
\[
  \vec{A} + \vec{B} = (A_x + B_x)\,\hat{i} + (A_y + B_y)\,\hat{j}
\]
\[
  |\vec{A}| = \sqrt{A_x^2 + A_y^2}, \qquad
  \theta = \arctan\!\left(\frac{A_y}{A_x}\right)
\]

\begin{pikabox}
Os vetores unitários $\hat{i}$, $\hat{j}$, $\hat{k}$ têm módulo 1 e apontam nas direções $x$, $y$ e $z$, respectivamente. São a \emph{base} do espaço cartesiano!
\end{pikabox}

\subsection{Coordenadas Polares}

Em trajetórias circulares ou curvas, é mais eficiente utilizar o sistema $(r, \theta)$:
\begin{itemize}[nosep]
  \item $x = r\cos\theta$
  \item $y = r\sin\theta$
\end{itemize}

\subsection{Produto Escalar e Produto Vetorial}

\[
  \vec{A} \cdot \vec{B} = A_x B_x + A_y B_y = |\vec{A}|\,|\vec{B}|\cos\theta
\]
\[
  \vec{A} \times \vec{B} = (A_y B_z - A_z B_y)\,\hat{i}
                          - (A_x B_z - A_z B_x)\,\hat{j}
                          + (A_x B_y - A_y B_x)\,\hat{k}
\]

\begin{charmanderbox}
$\vec{A} \cdot \vec{B}$ é \textbf{escalar}; $\vec{A} \times \vec{B}$ é \textbf{vetorial} e perpendicular ao plano de $\vec{A}$ e $\vec{B}$. Além disso, o produto vetorial \textbf{não é comutativo}: $\vec{A} \times \vec{B} = -\vec{B} \times \vec{A}$!
\end{charmanderbox}

% =====================================================================
\section{Velocidade e Aceleração Vetoriais}
% =====================================================================

\subsection{Vetor Deslocamento}

\[
  \Delta\vec{r} = \vec{r}_f - \vec{r}_i
\]

\begin{center}
\begin{tikzpicture}[>=stealth]
  \draw[very thin, gray!20] (-1,-1) grid (6,6);
  \draw[->, thick] (-1,0) -- (6,0) node[right] {$x$};
  \draw[->, thick] (0,-1) -- (0,6) node[above] {$y$};
  \coordinate (A) at (0,0);
  \coordinate (B) at (4,4);
  \filldraw[black] (A) circle (2pt) node[anchor=north east] {$P_i (0,0)$};
  \filldraw[black] (B) circle (2pt) node[anchor=south west] {$P_f (4,4)$};
  % Deslocamento vetorial
  \draw[->, dashed, gray, thick] (A) -- (B)
    node[midway, above left] {$\Delta \vec{r}$ (vetor)};
  % Trajetória real
  \draw[->, orange, ultra thick] (A) to [out=90, in=180] (2,3) to [out=0, in=225] (B);
  \node[orange, font=\bfseries] at (1.5,4) {Trajetória};
  \draw[dotted] (4,0) -- (4,4);
  \draw[dotted] (0,4) -- (4,4);
\end{tikzpicture}
\end{center}

\subsection{Velocidade Vetorial Média e Instantânea}

\[
  \bar{\vec{v}} = \frac{\Delta\vec{r}}{\Delta t}, \qquad
  \vec{v}(t) = \frac{d\vec{r}}{dt} = \dot{x}\,\hat{i} + \dot{y}\,\hat{j}
\]

A \textbf{velocidade escalar média} $v_{em}$ é a razão entre a distância total percorrida ($d$) e o tempo total, independentemente da direção:
\[
  v_{em} = \frac{d}{\Delta t}
\]

\begin{squirtlebox}
A \textbf{rapidez} (speed) é o módulo da velocidade:
\[
  |\vec{v}| = \sqrt{v_x^2 + v_y^2}
\]
A direção da velocidade é \textbf{sempre tangente} à trajetória!
\end{squirtlebox}

\subsection{Aceleração Vetorial Média e Instantânea}

\[
  \bar{\vec{a}} = \frac{\Delta\vec{v}}{\Delta t}, \qquad
  \vec{a}(t) = \frac{d\vec{v}}{dt} = \ddot{x}\,\hat{i} + \ddot{y}\,\hat{j}
\]

\begin{charmanderbox}
A aceleração pode ter \textbf{componente tangencial} (muda a rapidez) e \textbf{componente centrípeta} (muda a direção). Num movimento circular uniforme, só existe aceleração centrípeta!
\end{charmanderbox}

% =====================================================================
\section{Estudo de Caso I: A Jornada do Pikachu}
% =====================================================================

Nesta seção, analisamos o deslocamento do Pikachu do Laboratório do Prof.\ Carvalho até a padaria onde se encontra o Ash, aplicando os conceitos de cinemática vetorial.

\subsection{Ponto de Partida — Vetor Posição}

Definimos o Laboratório como a origem do sistema de coordenadas cartesianas.
\begin{itemize}[nosep]
  \item \textbf{Vetor Posição Inicial:} $\vec{r}_0 = 0\hat{i} + 0\hat{j}$
\end{itemize}

\subsection{O Desvio pelo Gramado — Coordenadas Polares}

Para evitar um bando de Spearows, o Pikachu realiza um movimento circular uniforme em um arco de $90^\circ$ ($\frac{\pi}{2}\,\text{rad}$), com raio $r = 10\,\text{m}$. Convertendo para coordenadas cartesianas:
\[
  \vec{r}_1 = (10\cos 90^\circ)\,\hat{i} + (10\sin 90^\circ)\,\hat{j} = 0\,\hat{i} + 10\,\hat{j} \;\text{[m]}
\]

\subsection{A Reta Final — Vetor Velocidade Média}

Pikachu avista a padaria em $\vec{r}_{\text{ash}} = 20\,\hat{i} + 10\,\hat{j}$ e percorre o trecho em $\Delta t = 5\,\text{s}$:
\begin{itemize}[nosep]
  \item \textbf{Deslocamento:} $\Delta\vec{r} = (20-0)\,\hat{i} + (10-10)\,\hat{j} = 20\,\hat{i}\;\text{m}$
  \item \textbf{Vetor Velocidade Média:}
\[
  \vec{v}_m = \frac{\Delta\vec{r}}{\Delta t} = \frac{20\,\hat{i}}{5} = 4\,\hat{i}\;\text{[m/s]}
\]
\end{itemize}

\subsection{Análise da Velocidade Escalar Média}

A distância total percorrida é a soma do comprimento do arco ($s$) com o trecho retilíneo final:
\[
  s = \tfrac{1}{4}(2\pi r) = \tfrac{1}{4}(2\pi \cdot 10) \approx 15{,}7\,\text{m}
\]
\[
  d = 15{,}7 + 20 = 35{,}7\,\text{m}
\]
Considerando um tempo total de $10\,\text{s}$:
\[
  v_{em} = \frac{d}{\Delta t_{\text{total}}} = \frac{35{,}7}{10} = 3{,}57\;\text{m/s}
\]

\begin{bulbabox}
O \textbf{vetor velocidade média} e a \textbf{velocidade escalar média} são diferentes! A velocidade média vetorial usa o deslocamento $\Delta\vec{r}$ (reta ponta-a-ponta), enquanto a escalar usa a distância total percorrida ao longo do caminho.
\end{bulbabox}

% =====================================================================
\section{Estudo de Caso II: A Jornada de Misty à Padaria}
% =====================================================================

Misty se encontra no ponto $A$ (origem) e precisa chegar à padaria no ponto $C$. A cidade tem traçado retangular, permitindo três trajetórias possíveis.

\subsection{Definição do Cenário}

\begin{itemize}[nosep]
  \item Ponto $A$ (Misty): $(0, 0)$
  \item Ponto $C$ (Padaria): $(4, 3)$
  \item Distância horizontal: $a = 4$ unidades; distância vertical: $b = 3$ unidades.
\end{itemize}

\begin{center}
\begin{tikzpicture}[>=stealth, scale=1.2]
  \draw[very thin, gray!20] (-1,-1) grid (5,4);
  \draw[->, thick] (-0.5,0) -- (5,0) node[right] {$x$};
  \draw[->, thick] (0,-0.5) -- (0,4) node[above] {$y$};
  \coordinate (A) at (0,0);
  \coordinate (B) at (4,0);
  \coordinate (C) at (4,3);
  \coordinate (D) at (0,3);
  \filldraw[black] (A) circle (2pt) node[anchor=north east, font=\bfseries] {A (Misty)};
  \filldraw[black] (B) circle (2pt) node[anchor=north west] {B};
  \filldraw[black] (C) circle (2pt) node[anchor=south west, font=\bfseries] {C (Padaria)};
  \filldraw[black] (D) circle (2pt) node[anchor=south east] {D};
  % Trajetória 1: ABC
  \draw[->, thick, pokeblue] (A) -- (B) node[midway, below] {$a$};
  \draw[->, thick, pokeblue] (B) -- (C) node[midway, right] {$b$};
  \node[pokeblue] at (2,-0.4) {$\vec{r}_1$};
  % Trajetória 2: ADC
  \draw[->, thick, pokered, dashed] (A) -- (D) node[midway, left] {$b$};
  \draw[->, thick, pokered, dashed] (D) -- (C) node[midway, above] {$a$};
  \node[pokered] at (-0.4, 1.5) {$\vec{r}_2$};
  % Trajetória 3: atalho
  \draw[->, ultra thick, pokegreen] (A) -- (C);
  \node[pokegreen, rotate=36.8] at (2, 1.8) {$\vec{r}_3$ (Atalho)};
  \node[anchor=north west, font=\footnotesize] at (4.2, 1.5) {$a=4,\;b=3$};
\end{tikzpicture}
\end{center}

\subsection{Análise das Trajetórias}

O \textbf{vetor deslocamento} é sempre $\Delta\vec{r} = 4\,\hat{i} + 3\,\hat{j}$ para os três caminhos. O que difere é a \textbf{distância percorrida}:

\subsubsection{Trajetória $\vec{r}_1$ — Caminho $ABC$}
\[
  d_1 = a + b = 4 + 3 = 7\;\text{unidades}
\]

\subsubsection{Trajetória $\vec{r}_2$ — Caminho $ADC$}
\[
  d_2 = b + a = 3 + 4 = 7\;\text{unidades}
\]

\subsubsection{Trajetória $\vec{r}_3$ — Atalho $AC$}
\[
  d_3 = \sqrt{a^2 + b^2} = \sqrt{4^2 + 3^2} = \sqrt{25} = 5\;\text{unidades}
\]

\begin{squirtlebox}
\textbf{Conclusão:} O \textbf{vetor deslocamento} ($\Delta\vec{r} = 4\hat{i}+3\hat{j}$) foi \textit{idêntico} para as três trajetórias, pois depende apenas das posições inicial e final. A \textbf{distância percorrida} variou: $7$ nas trajetórias urbanas e $5$ no atalho. Misty escolhe $\vec{r}_3$ e chega mais rápido à padaria!
\end{squirtlebox}

% =====================================================================
\section{Análise Estrutural do Deslocamento Vetorial}
% =====================================================================

\subsection{Definição do Vetor Deslocamento Parametrizado}

Consideramos um deslocamento vetorial $\Delta\vec{r}$ definido pelas componentes de escala $R$ e $T$:
\begin{equation}
  \Delta\vec{r} = -RT\,\hat{i} + RT\,\hat{j}
\end{equation}

\subsection{Módulo do Deslocamento}

\[
  |\Delta\vec{r}| = \sqrt{(-RT)^2 + (RT)^2} = \sqrt{2R^2T^2} = RT\sqrt{2}
\]

\subsection{Distância Escalar pela Trajetória Ortogonal}

Se o deslocamento ocorrer de forma ortogonal (seguindo primeiro o eixo $x$ e depois o eixo $y$):
\[
  \Delta s = |-RT| + |RT| = 2RT
\]

\subsection{Comparação entre as Grandezas}

\begin{center}
\renewcommand{\arraystretch}{1.4}
\begin{tabular}{|l|c|}
\hline
\textbf{\color{pokered}Grandeza} & \textbf{\color{pokered}Valor} \\
\hline
\color{pokeblue}Deslocamento vetorial (módulo) & $RT\sqrt{2} \approx 1{,}41\,RT$ \\
\hline
\color{pokeblue}Distância escalar (trajetória retangular) & $2RT$ \\
\hline
\end{tabular}
\end{center}

Como $2RT > 1{,}41RT$, o caminho percorrido ao longo dos eixos é aproximadamente $41\%$ mais longo do que o deslocamento direto entre os dois pontos.

\begin{center}
\begin{tikzpicture}[>=stealth]
  \draw[->, thin, gray] (-3,0) -- (1,0) node[right] {$x$};
  \draw[->, thin, gray] (0,-1) -- (0,3) node[above] {$y$};
  \draw[->, pokeblue, ultra thick] (0,0) -- (-2,2)
    node[midway, above left] {$\Delta\vec{r}$};
  \draw[dashed, pokered] (0,0) -- (-2,0) -- (-2,2);
  \node[pokered, below] at (-1,0) {$-RT\hat{i}$};
  \node[pokered, left] at (-2,1) {$RT\hat{j}$};
\end{tikzpicture}
\end{center}

% =====================================================================
\section{Análise Curvilínea: Trajetória Circular e Comparação}
% =====================================================================

\subsection{Deslocamento Vetorial}

Mantemos $\Delta\vec{r} = -RT\,\hat{i} + RT\,\hat{j}$, cujo módulo (a corda do arco) é:
\[
  |\Delta\vec{r}| = RT\sqrt{2} \approx 1{,}41\,RT
\]

\subsection{Distância Escalar pelo Arco Circular}

Para um arco de um quarto de circunferência ($90^\circ = \frac{\pi}{2}\,\text{rad}$) com raio $RT$:
\[
  \Delta s = \frac{\pi}{2}\cdot RT \approx 1{,}57\,RT
\]

\subsection{Comparação Analítica}

\begin{itemize}
  \item \textbf{Caminho mais curto (vetorial):} $1{,}41\,RT$
  \item \textbf{Caminho curvo (escalar):} $1{,}57\,RT$
\end{itemize}

A razão entre as distâncias é $\dfrac{\pi}{2\sqrt{2}} \approx 1{,}11$: o arco circular é aproximadamente $11\%$ mais longo que o deslocamento direto.

\begin{center}
\begin{tikzpicture}[>=stealth]
  \draw[->, gray!50] (-3,0) -- (1,0) node[right] {$x$};
  \draw[->, gray!50] (0,-1) -- (0,3) node[above] {$y$};
  % Vetor deslocamento
  \draw[->, pokeblue, thick] (0,0) -- (-2,2)
    node[midway, below left] {$|\Delta\vec{r}|$};
  % Arco circular
  \draw[pokeorange, ultra thick] (0,0) arc (0:90:-2);
  \node[pokeorange] at (-1.8,1.8) {$\Delta s = \frac{\pi RT}{2}$};
  \filldraw (0,0) circle (1.5pt) node[below right] {$P_i$};
  \filldraw (-2,2) circle (1.5pt) node[above left] {$P_f$};
\end{tikzpicture}
\end{center}

% =====================================================================
\section{Movimento 2D com Aceleração Uniforme}
% =====================================================================

\subsection{Equações Cinemáticas Vetoriais}

Para $\vec{a} = $ constante:

\begin{definicaobox}
\begin{align}
  \vec{v}(t) &= \vec{v}_0 + \vec{a}\,t \\[4pt]
  \vec{r}(t) &= \vec{r}_0 + \vec{v}_0\,t + \tfrac{1}{2}\vec{a}\,t^2 \\[4pt]
  v^2 &= v_0^2 + 2\vec{a}\cdot(\vec{r}-\vec{r}_0)
\end{align}
\end{definicaobox}

\subsection{Lançamento de Projéteis}

Tomamos $\vec{a} = -g\,\hat{j}$ ($g \approx 9{,}8\,\text{m/s}^2$):
\[
  x(t) = x_0 + v_{0x}\,t, \qquad y(t) = y_0 + v_{0y}\,t - \tfrac{1}{2}g\,t^2
\]
\[
  v_x = v_0\cos\theta_0, \qquad v_y = v_0\sin\theta_0 - g\,t
\]

\textbf{Alcance máximo} (para $y_0 = 0$):
\[
  R = \frac{v_0^2 \sin 2\theta_0}{g}
\]
\textbf{Altura máxima:}
\[
  H = \frac{(v_0 \sin\theta_0)^2}{2g}
\]

\begin{pikabox}
O alcance é máximo quando $\theta_0 = 45°$. Ângulos complementares ($\theta$ e $90° - \theta$) produzem o mesmo alcance!
\end{pikabox}

\begin{bulbabox}
\textbf{Exemplo:} Um projétil é lançado com $v_0 = 30\,\text{m/s}$ a $\theta_0 = 60°$.
\[
  v_{0x} = 30\cos 60° = 15\,\text{m/s}, \quad v_{0y} = 30\sin 60° = 15\sqrt{3}\,\text{m/s}
\]
\[
  H = \frac{(15\sqrt{3})^2}{2 \times 9{,}8} = \frac{675}{19{,}6} \approx 34{,}4\,\text{m}
\]
\[
  R = \frac{(30)^2 \sin 120°}{9{,}8} = \frac{900 \times \frac{\sqrt{3}}{2}}{9{,}8} \approx 79{,}5\,\text{m}
\]
\end{bulbabox}

\subsection{Movimento Circular Uniforme (MCU)}

\[
  |\vec{a}_c| = \frac{v^2}{R}, \quad \vec{a}_c \perp \vec{v}
\]
\[
  T = \frac{2\pi R}{v}, \quad f = \frac{1}{T}, \quad \omega = 2\pi f = \frac{v}{R}
\]

\begin{squirtlebox}
No MCU, embora a \textbf{rapidez} seja constante, a \textbf{velocidade vetorial} muda de direção continuamente. Isso exige uma aceleração centrípeta ($a_c$) sempre voltada ao centro da trajetória.
\end{squirtlebox}

% =====================================================================
\section{Velocidade Relativa e Composição de Movimentos}
% =====================================================================

\subsection{Transformação de Galileu}

Consideramos dois referenciais: o solo (Misty, fixo) e o referencial em movimento (Ash). A velocidade do Pikachu em relação ao solo é:
\begin{equation}
  \vec{v}_{P,S} = \vec{v}_{P,A} + \vec{v}_{A,S}
\end{equation}

A velocidade relativa observada pelo Ash é:
\begin{equation}
  \vec{v}_{P,A} = \vec{v}_{P,S} - \vec{v}_{A,S}
\end{equation}

\subsection{Representação Gráfica: Triângulo de Velocidades}

\begin{center}
\begin{tikzpicture}[>=stealth, scale=0.8]
  % Gráfico 1: Posição
  \begin{scope}[xshift=-5cm]
    \draw[->, thin, gray] (-1,0) -- (4,0) node[right] {$x$};
    \draw[->, thin, gray] (0,-1) -- (0,4) node[above] {$y$};
    \coordinate (M) at (0,0);
    \coordinate (A) at (2,1);
    \coordinate (P) at (3,3);
    \draw[->, thick, pokeblue] (M) -- (A) node[midway, below right] {$\vec{r}_{A,S}$};
    \draw[->, thick, pokered] (A) -- (P) node[midway, left] {$\vec{r}_{P,A}$};
    \draw[->, ultra thick, pokeorange] (M) -- (P) node[midway, above left] {$\vec{r}_{P,S}$};
    \filldraw (M) circle (2pt) node[below left] {Misty (S)};
    \filldraw (A) circle (2pt) node[right] {Ash (S')};
    \filldraw (P) circle (2pt) node[above] {Pikachu};
    \node at (1.5,-1.5) {(a) Vetores Posição};
  \end{scope}
  % Gráfico 2: Triângulo de velocidades
  \begin{scope}[xshift=4cm]
    \draw[->, thin, gray] (-1,0) -- (4,0) node[right] {$x$};
    \draw[->, thin, gray] (0,-1) -- (0,4) node[above] {$y$};
    \draw[->, ultra thick, pokeblue] (0,0) -- (3,0) node[midway, below] {$\vec{v}_{A,S}$};
    \draw[->, ultra thick, pokered] (3,0) -- (3,2) node[midway, right] {$\vec{v}_{P,A}$};
    \draw[->, ultra thick, pokeorange] (0,0) -- (3,2) node[midway, above left] {$\vec{v}_{P,S}$};
    \node at (1.5,-1.5) {(b) Triângulo de Velocidades};
  \end{scope}
\end{tikzpicture}
\end{center}

\subsection{Análise do Módulo}

Quando os movimentos são perpendiculares, o módulo da velocidade resultante é obtido pelo Teorema de Pitágoras:
\begin{equation}
  v_{P,S} = \sqrt{v_{A,S}^2 + v_{P,A}^2}
\end{equation}

\begin{bulbabox}
\textbf{Exemplo:} Ash se move a $4\,\text{m/s}$ para o leste e o Pikachu se move a $3\,\text{m/s}$ para o norte em relação ao Ash. A velocidade que Misty observa é:
\[
  v_{P,S} = \sqrt{4^2 + 3^2} = \sqrt{25} = 5\,\text{m/s}
\]
\end{bulbabox}

% =====================================================================
\section{Dedução Analítica via Cálculo Diferencial}
% =====================================================================

\subsection{Obtenção da Velocidade Relativa}

Seja $\vec{r}_{PS}$ a posição do Pikachu em relação ao solo, $\vec{r}_{PM}$ em relação ao Ash, e $\vec{r}_{MS}$ a posição do Ash no solo:
\[
  \vec{r}_{PS} = \vec{r}_{PM} + \vec{r}_{MS}
\]

Derivando em relação ao tempo:
\[
  \frac{d\vec{r}_{PS}}{dt} = \frac{d\vec{r}_{PM}}{dt} + \frac{d\vec{r}_{MS}}{dt}
  \quad\Rightarrow\quad
  \vec{v}_{PS} = \vec{v}_{PM} + \vec{v}_{MS}
\]

\subsection{Análise da Aceleração e Referenciais Inerciais}

Como o Ash se move com velocidade constante ($\vec{v}_{MS} = \vec{v} = \text{cte}$):
\[
  \vec{a}_{MS} = \frac{d\vec{v}_{MS}}{dt} = 0
\]

Derivando a equação de velocidade:
\[
  \vec{a}_{PS} = \vec{a}_{PM}
\]

\begin{squirtlebox}
O fato de $\dfrac{d\vec{v}_{MS}}{dt} = 0$ implica que o referencial do Ash é um \textbf{Referencial Inercial}. Isso significa que a aceleração observada por Misty e pelo Ash é a mesma — as leis de Newton permanecem idênticas para ambos os observadores!
\end{squirtlebox}

\begin{center}
\begin{tikzpicture}[>=stealth, scale=1.1]
  \draw[->, pokeblue, ultra thick] (0,0) -- (3,0)
    node[midway, below] {$\vec{v}_{MS} = \vec{v}$};
  \draw[->, pokered, ultra thick] (3,0) -- (3,2)
    node[midway, right] {$\vec{v}_{PM}$};
  \draw[->, pokeorange, ultra thick] (0,0) -- (3,2)
    node[midway, above left] {$\vec{v}_{PS}$};
  \node[anchor=north] at (1.5,-0.5) {\footnotesize $\vec{v}_{PS} = \vec{v}_{PM} + \vec{v}_{MS}$};
\end{tikzpicture}
\end{center}

% =====================================================================
\section{Dinâmica: As Leis de Newton}
% =====================================================================

\begin{definicaobox}
\textbf{2ª Lei (Princípio Fundamental):} A força resultante é o produto da massa pela aceleração:
\[
  \sum\vec{F} = m\,\vec{a}
\]
\end{definicaobox}

\begin{pikabox}
\textbf{Machamp explica:} Se você empurra um Snorlax com uma força $F$, ele terá uma aceleração pequena devido à sua grande massa (inércia). Já um Rattata voaria longe com a mesma força!
\end{pikabox}

% =====================================================================
\section{Trabalho e Energia}
% =====================================================================

\begin{squirtlebox}
O Trabalho ($W$) é a transferência de energia através de uma força:
\[
  W = F \cdot d \cdot \cos\theta
\]
Pelo Teorema Trabalho--Energia Cinética:
\[
  W_{\text{res}} = \Delta K = \tfrac{1}{2}mv_f^2 - \tfrac{1}{2}mv_i^2
\]
\end{squirtlebox}

\begin{bulbabox}
\textbf{Snorlax em Repouso:} Para tirar um Snorlax ($m=460\,\text{kg}$) da inércia e fazê-lo atingir $2\,\text{m/s}$, quanto trabalho é necessário?
\[
  W = \tfrac{1}{2}(460)(2^2) = 920\,\text{Joules}
\]
\end{bulbabox}

% =====================================================================
\section{Resumo das Fórmulas Essenciais}
% =====================================================================

\begin{center}
\renewcommand{\arraystretch}{1.6}
\arrayrulecolor{pokeblue}
\begin{tabular}{|l|c|}
\hline
\rowcolor{pokered} \textbf{\color{white}Grandeza} & \textbf{\color{white}Fórmula} \\
\hline
\color{pokeblue}Velocidade média          & $\bar{v} = \Delta x / \Delta t$ \\
\hline
\color{pokeblue}Velocidade instantânea    & $v = dx/dt$ \\
\hline
\color{pokeblue}Aceleração média          & $\bar{a} = \Delta v / \Delta t$ \\
\hline
\color{pokeblue}MRU (posição)             & $x = x_0 + vt$ \\
\hline
\color{pokeblue}MRUV (velocidade)         & $v = v_0 + at$ \\
\hline
\color{pokeblue}MRUV (posição)            & $x = x_0 + v_0 t + \tfrac{1}{2}at^2$ \\
\hline
\color{pokeblue}Torricelli                & $v^2 = v_0^2 + 2a\,\Delta x$ \\
\hline
\color{pokeblue}Módulo de vetor 2D        & $|\vec{A}| = \sqrt{A_x^2+A_y^2}$ \\
\hline
\color{pokeblue}Vel.\ vetorial instantânea& $\vec{v} = d\vec{r}/dt$ \\
\hline
\color{pokeblue}Aceleração vetorial       & $\vec{a} = d\vec{v}/dt$ \\
\hline
\color{pokeblue}Alcance (projétil)        & $R = v_0^2 \sin 2\theta_0 / g$ \\
\hline
\color{pokeblue}Altura máxima             & $H = (v_0\sin\theta_0)^2 / (2g)$ \\
\hline
\color{pokeblue}Aceleração centrípeta     & $a_c = v^2/R$ \\
\hline
\color{pokeblue}Período (MCU)             & $T = 2\pi R / v$ \\
\hline
\color{pokeblue}Velocidade angular        & $\omega = 2\pi f = v/R$ \\
\hline
\color{pokeblue}Força resultante          & $\sum\vec{F} = m\,\vec{a}$ \\
\hline
\color{pokeblue}Energia cinética          & $K = \tfrac{1}{2}mv^2$ \\
\hline
\color{pokeblue}Trabalho                  & $W = F\,d\,\cos\theta$ \\
\hline
\color{pokeblue}Vel.\ relativa (Galileu)  & $\vec{v}_{PS} = \vec{v}_{PM} + \vec{v}_{MS}$ \\
\hline
\end{tabular}
\end{center}

\vfill
\begin{center}
\textit{``Gotta catch 'em all --- e entender todos os movimentos!''}\\[4pt]
\textcolor{pokered}{\rule{0.5\textwidth}{1pt}}\\[4pt]
\small Apostila elaborada com base nas aulas de Fundamentos de Mecânica --- UAI Física.\\
\small Profª Camila Leite Coura Mariano de Oliveira.
\end{center}

\end{document}
